\documentclass[11pt,a4paper,oneside]{report}             % Single-side
%\documentclass[11pt,a4paper,twoside,openright]{report}  % Duplex

%\PassOptionsToPackage{chapternumber=Huordinal}{magyar.ldf}
\usepackage{t1enc}
\usepackage[utf8]{inputenc}
\usepackage{amsmath}
\usepackage{amssymb}
\usepackage[english,magyar]{babel}
\usepackage{datetime}
\usepackage{enumerate}
\usepackage[thmmarks]{ntheorem}
\usepackage{graphics}
\usepackage{epsfig}
\usepackage{listings}
\usepackage{color}
%\usepackage{fancyhdr}
\usepackage{lastpage}
\usepackage{anysize}
\usepackage{sectsty}
\usepackage{setspace}  % Ettol a tablazatok, abrak, labjegyzetek maradnak 1-es sorkozzel!
\usepackage[hang]{caption}
\usepackage{hyperref}

%--------------------------------------------------------------------------------------
% Main variables
%--------------------------------------------------------------------------------------
\newcommand{\vikszerzo}{Dudás Ádám}
\newcommand{\vikkonzulensa}{Bergmann Gábor}
\newcommand{\vikkonzulensb}{Ujhelyi Zoltán}
\newcommand{\vikcim}{Modell-lekérdezések statikus analízise}
\newcommand{\viktanszek}{Méréstechnika és Információs Rendszerek Tanszék}
\newcommand{\vikdoktipus}{Szakdolgozat}

%--------------------------------------------------------------------------------------
% Page layout setup
%--------------------------------------------------------------------------------------
% we need to redefine the pagestyle plain
% another possibility is to use the body of this command without \fancypagestyle
% and use \pagestyle{fancy} but in that case the special pages
% (like the ToC, the References, and the Chapter pages)remain in plane style

\pagestyle{plain}
%\setlength{\parindent}{0pt} % áttekinthetőbb, angol nyelvű dokumentumokban jellemző
%\setlength{\parskip}{8pt plus 3pt minus 3pt} % áttekinthetőbb, angol nyelvű dokumentumokban jellemző
\setlength{\parindent}{12pt} % magyar nyelvű dokumentumokban jellemző
\setlength{\parskip}{0pt}    % magyar nyelvű dokumentumokban jellemző

\marginsize{35mm}{25mm}{15mm}{15mm} % anysize package
\setcounter{secnumdepth}{0}
\sectionfont{\large\upshape\bfseries}
\setcounter{secnumdepth}{2}
\singlespacing
\frenchspacing

%---------------------------------------------------------------------------
% This makes Chapter go away, so
%   Chapter 1
%   My first chapter
% becomes
%   1. My first chapter
%
\makeatletter
\renewcommand{\@makechapterhead}[1]{%
\vspace*{50 pt}%
{\setlength{\parindent}{0pt} \raggedright \normalfont
\bfseries\Huge\thechapter.\ #1
\par\nobreak\vspace{40 pt}}}
\makeatother
%---------------------------------------------------------------------------

%--------------------------------------------------------------------------------------
%	Setup hyperref package
%--------------------------------------------------------------------------------------
\hypersetup{
    bookmarks=true,            % show bookmarks bar?
    unicode=false,             % non-Latin characters in Acrobat’s bookmarks
    pdftitle={\vikcim},        % title
    pdfauthor={\vikszerzo},    % author
    pdfsubject={\vikdoktipus}, % subject of the document
    pdfcreator={\vikszerzo},   % creator of the document
    pdfproducer={Producer},    % producer of the document
    pdfkeywords={keywords},    % list of keywords
    pdfnewwindow=true,         % links in new window
    colorlinks=true,           % false: boxed links; true: colored links
    linkcolor=black,           % color of internal links
    citecolor=black,           % color of links to bibliography
    filecolor=black,           % color of file links
    urlcolor=black             % color of external links
}

%--------------------------------------------------------------------------------------
% Set up listings
%--------------------------------------------------------------------------------------
\lstset{
	basicstyle=\scriptsize\ttfamily, % print whole listing small
	keywordstyle=\color{black}\bfseries\underbar, % underlined bold black keywords
	identifierstyle=, 					% nothing happens
	commentstyle=\color{white}, % white comments
	stringstyle=\scriptsize\sffamily, 			% typewriter type for strings
	showstringspaces=false,     % no special string spaces
	aboveskip=3pt,
	belowskip=3pt,
	columns=fixed,
	backgroundcolor=\color{lightgray},
	extendedchars=true,
	literate={á}{{\'a}}1 {é}{{\'e}}1 {í}{{\'i}}1 {ó}{{\'o}}1 {ö}{{\"o}}1 {ő}{{\H{o}}}1 {ú}{{\'u}}1 {ü}{{\"u}}1 {ű}{{\H{u}}}1,
} 		
\def\lstlistingname{lista}	

%--------------------------------------------------------------------------------------
%	Some new commands and declarations
%--------------------------------------------------------------------------------------
\newcommand{\code}[1]{{\upshape\ttfamily\scriptsize\indent #1}}

% define references
\newcommand{\figref}[1]{\ref{fig:#1}.}
\renewcommand{\eqref}[1]{(\ref{eq:#1})}
\newcommand{\listref}[1]{\ref{listing:#1}.}
\newcommand{\sectref}[1]{\ref{sect:#1}}
\newcommand{\tabref}[1]{\ref{tab:#1}.}

\DeclareMathOperator*{\argmax}{arg\,max}
%\DeclareMathOperator*[1]{\floor}{arg\,max}
\DeclareMathOperator{\sign}{sgn}
\DeclareMathOperator{\rot}{rot}
\definecolor{lightgray}{rgb}{0.95,0.95,0.95}

\author{\vikszerzo}
\title{\vikcim}
\includeonly{
	titlepage,%
	declaration,%
	abstract,%
	introduction,%
	chapter1,%
	chapter2,%
	chapter3,%
	acknowledgement,%
	appendices,%
}

%--------------------------------------------------------------------------------------
%	Setup captions
%--------------------------------------------------------------------------------------
\captionsetup[figure]{
%labelsep=none,
%font={footnotesize,it},
%justification=justified,
width=.75\textwidth,
aboveskip=10pt}

\renewcommand{\captionlabelfont}{\small\bf}
\renewcommand{\captionfont}{\footnotesize\it}

%--------------------------------------------------------------------------------------
% Table of contents and the main text
%--------------------------------------------------------------------------------------
\begin{document}
\singlespacing
%\include{guideline}
%\include{project}

\pagenumbering{arabic}
\onehalfspacing
%--------------------------------------------------------------------------------------
%	The title page
%--------------------------------------------------------------------------------------
\begin{titlepage}
\begin{center}
\includegraphics[width=60mm,keepaspectratio]{figures/BMElogo.png}\\
\vspace{0.3cm}
\textbf{Budapesti Műszaki és Gazdaságtudományi Egyetem}\\
\textmd{Villamosmérnöki és Informatikai Kar}\\
\textmd{\viktanszek}\\[5cm]

\vspace{0.4cm}
{\huge \bfseries \vikcim}\\[0.8cm]
\vspace{0.5cm}
\textsc{\Large \vikdoktipus}\\[4cm]

\begin{tabular}{cc}
 \makebox[7cm]{\emph{Készítette}} & \makebox[7cm]{\emph{Konzulensek}} \\
 \makebox[7cm]{\vikszerzo} & \makebox[7cm]{\vikkonzulensa} \\
                           & \makebox[7cm]{\vikkonzulensb}
\end{tabular}

\vfill
{\large \today}
\end{center}
\end{titlepage}



\tableofcontents\vfill
%--------------------------------------------------------------------------------------
% Nyilatkozat
%--------------------------------------------------------------------------------------
\begin{center}
\large
\textbf{HALLGATÓI NYILATKOZAT}\\
\end{center}

Alulírott \emph{\vikszerzo}, szigorló hallgató kijelentem, hogy ezt a szakdolgozatot meg nem engedett segítség nélkül, saját magam készítettem, csak a megadott forrásokat (szakirodalom, eszközök stb.) használtam fel. Minden olyan részt, melyet szó szerint, vagy azonos értelemben, de átfogalmazva más forrásból átvettem, egyértelműen, a forrás megadásával megjelöltem.

Hozzájárulok, hogy a jelen munkám alapadatait (szerző(k), cím, angol és magyar nyelvű tartalmi kivonat, készítés éve, konzulens(ek) neve) a BME VIK nyilvánosan hozzáférhető elektronikus formában, a munka teljes szövegét pedig az egyetem belső hálózatán keresztül (vagy autentikált felhasználók számára) közzétegye. Kijelentem, hogy a benyújtott munka és annak elektronikus verziója megegyezik. Dékáni engedéllyel titkosított diplomatervek esetén a dolgozat szövege csak 3 év eltelte után válik hozzáférhetővé.

\begin{flushleft}
\vspace*{1cm}
Budapest, \today
\end{flushleft}

\begin{flushright}
 \vspace*{1cm}
 \makebox[7cm]{\rule{6cm}{.4pt}}\\
 \makebox[7cm]{\emph{\vikszerzo}}\\
 \makebox[7cm]{hallgató}
\end{flushright}
\thispagestyle{empty}

\vfill
\clearpage
\thispagestyle{empty} % an empty page


%----------------------------------------------------------------------------
% Abstract in hungarian
%----------------------------------------------------------------------------
\chapter*{Kivonat}\addcontentsline{toc}{chapter}{Kivonat}

A \emph{modell-vezérelt szoftverfejlesztés} egy manapság széles körben elterjedt technika, mely megkönnyíti mind a szoftverek előállítását, mind azok későbbi továbbfejlesztését, karbantartását.
A technika alkalmazása során a fejlesztők az adott probléma doménjének egy olyan modelljét, modelljeit készítik el, mely(ek) a probléma megoldásának szempontjából fontos részeit, így segítve a megismerést, a megértést és a szimulációt.
\vfill

%----------------------------------------------------------------------------
% Abstract in english
%----------------------------------------------------------------------------
\chapter*{Abstract}\addcontentsline{toc}{chapter}{Abstract}

This document is a \LaTeX-based skeleton for BSc/MSc~theses of students at the Electrical Engineering and Informatics Faculty, Budapest University of Technology and Economics. The usage of this skeleton is optional. It has been tested with the \emph{TeXLive} \TeX~implementation, and it requires the PDF-\LaTeX~compiler.
\vfill


%----------------------------------------------------------------------------
\chapter{Bevezető}%\addcontentsline{toc}{chapter}{Bevezető}
%----------------------------------------------------------------------------

% A bevezető tartalmazza a diplomaterv-kiírás elemzését, történelmi előzményeit, a feladat indokoltságát (a motiváció leírását), az eddigi megoldásokat, és ennek tükrében a hallgató megoldásának összefoglalását.

Napjainkra a szoftverek összetettsége jelentősen megnőtt az egyre fejlettebb hardverek adta lehetőségek és a felhasználók olykor végeláthatatlan igényeinek következtében.
Ezen komplexitás velejárója, hogy a szoftverek tervezése, készítése, és néha már a használatuk is komoly kihívást jelenthet fejlesztőik és felhasználóik számára.
Az igények kielégítése mellett fontos a tervezési és a készítési folyamat, illetve a későbbi továbbfejlesztési és szinte elkerülhetetlen karbantartási munkák idő- és erőforrásigényének minimalizálása.
Ráadásul nem árt, ha mindezek mellett szoftverünk jól dokumentált a felhasználók és a fejlesztők számára egyaránt.  

\section{Modellezés}

A \emph{modell} definíciója szerint egy rendszer, elmélet, vagy jelenség leírása, mely rendelkezik annak ismert, vagy kikövetkeztetett tulajdonságaival és a jellemzői további vizsgálatára használható \cite{dict:Model}.
A mérnöki diszciplínákban szinte mindenhol alkalmaznak matematikai vagy egyéb modelleket a komplex problémák leírására, vizualizálására.
Ezek a modellek azáltal segítik a mérnökök munkáját, hogy a -- valós vagy elképzelt -- világnak csak a problémához szorosan kapcsolódó, a megoldás szempontjából lényeges részeit reprezentálják -- írják le vagy jelenítik meg.
A modellek manipulációján keresztül a mérnököknek lehetősége adódik ötleteiket kipróbálni, terveiket letesztelni különösebb erőforrás ráfordítás nélkül.

A modellezés tehát olyan eszközt biztosít számunkra, mellyel egy adott rendszerről következtetéseket vonhatunk le.
Felmerülhet a kérdés, hogy vajon milyen eszközzel tudjuk magukat a modellezéshez használt eszközöket vizsgálni?
A válasz szinte triviális: a modellező eszközöket is tudjuk modellezni.
Ez a gondolat vezet el a \emph{modellezési szintek} fogalmához.
A dolgozat kapcsán az alábbi három modellezési szintet érdemes megkülönböztetnünk:
\begin{itemize}
    \item a \emph{(példány-)modellek} szintje a legalacsonyabb, itt találhatóak a konkrét rendszereket reprezentáló modellek (példa: a $2 + 3i$ komplex szám két valós számból álló tömbként a számítógép memóriájában)
    \item eggyel magasabb absztrakciós szintet képvisel a példánymodellek struktúráját leíró \emph{metamodellek} szintje (példa: a komplex számokat leíró \gls{UML} osztálydiagram)
    \item és végül a metamodellek általános felépítését specifikáló \emph{meta-metamodell} szintről (példa: az \gls{UML} specifikációja)
\end{itemize}
Nyilvánvalóan a lista a magasabb absztrakciós szintek irányába tetszőleges hosszan folytatható, amíg az aktuális legmagasabb szinten lévő modelleket leíró közös struktúra meghatározható.
Továbbá az egyes szinteken található modellek nagyban függenek a modellezendő rendszer megválasztásától, emiatt ugyanaz a modell két különböző választás esetén két különböző szintet is képviselhet.

Nem meglepő módon, mint modern mérnöki tudomány, az informatika is alkalmaz modelleket.
A szoftverrendszerek tervezésénél régóta használnak többféle modellt:
\begin{itemize}
	\item adatszerkezeti modell (objektummodell, osztályhierarchia, öröklődés, stb.)
	\item funkcionális modell (adatforrások, -tárolók, -nyelők, -folyamok)
	\item dinamikus modell (állapotok, események, kommunikáció) \cite{Kondorosi07}.
\end{itemize}
Az ezen modelleken alapuló nézetek és diagramok -- osztályhierarchia, külső felhasználók, objektumok együttműködése, tevékenységek, események időbelisége, objektumok állapota, rendelkezésre álló szoftver- és hardver erőforrások -- nélkül a tervezés a néha akár több millió kódsorból álló szoftverrendszerek esetében áttekinthetetlen, a rendszer implementálás előtti tesztelése, viselkedésének elemzése, idő- és költségbecslések elvégzése kivitelezhetetlen volna.

A modellezés fejlődése a szoftverfejlesztésben a 2000-es évekre elvezetett a \emph{modell-vezérelt paradigmák} -- az általános model-driven software developmenttől (\gls{MDSD}), a speciálisabb model-driven engineeringen (\gls{MDE}) át, az \gls{OMG} által konkrétan specifikált model-driven architecture-ig (\gls{MDA}) -- és eszközrendszerek kifejlődéséhez \cite{OMG:MDA:FAQ,Schmidt:MDE:10.1109/MC.2006.58,MDSD:TEM06}.
A modell-vezérelt szoftverfejlesztés (\gls{MDSD}) egy olyan, manapság széles körben elterjedt \gls{CASE} technika, mely megkönnyíti mind a szoftverek előállítását, mind azok későbbi továbbfejlesztését, karbantartását \cite{MDSD:TEM06}.
A technika alkalmazása során a fejlesztők először az adott feladat alapján egy domén-specifikus, magas szintű, platformfüggetlen \emph{modellt hoznak létre}.
Ezután az elkészült modellt különböző szoftvereszközök segítségével \emph{transzformálják} más alakokba, például dokumentációvá, grafikonokká, programkóddá, végrehajtható programokká vagy egyéb modellekké.
Az ilyen szoftveres, automatizált transzformáció nagy előnye, hogy a fejlesztők által meghatározott módon képes a modell több kapcsolódó reprezentációja közötti -- bizonyos esetekben akár kétirányú -- konzisztencia megőrzésére.
Teszi mindezt hosszútávon kevesebb befektetett munkával, a kiküszöbölt hibalehetőségek számáról nem is beszélve.
Emellett az egyes transzformációk -- összeillő ki- és bemenetek mentén -- össze is kapcsolhatóak, így egyszerűbb transzformációkból egész transzformációs csővezetékeket (transformation pipeline) hozhatunk létre.
A transzformációk megadásához általában egy modelltranszformációs nyelvet (\gls{MTL}) használunk, amilyen pl. a VIATRA2 Textual Command Language \cite{Balogh:2006:AMT:1141277.1141575}.

A modelltranszformációk egy speciális esete a \emph{lekérdezés} (query), mely a modell reprezentációján nem változtat, viszont csak azon elemeit adja vissza, melyek kielégítenek egy -- akár többszörösen összetett -- feltételt.
Az ilyen modell-lekérdezéseket felhasználhatjuk többek között a modell inkonzisztenciáinak kimutatására, származtatott attribútumok és kapcsolatok implementálására, modelltranszformációk szabály alapú vezérlésére, vagy nézetek deklaratív definiálására.

Mint minden munkafolyamat, amelyben emberi közreműködésre van szükség, a modellek készítése és manipulálása, és a transzformációk leírása is temérdek lehetőséget ad az embernek hibák elkövetésére.
A mérnökök ezért megpróbálják az ilyen tevékenységek mennyiségét csökkenteni azáltal, hogy előre elkészített és fokozottan ellenőrzött komponenseket biztosítanak a fejlesztők számára függvénykönyvtárak vagy keretrendszerek formájában.
A fennmaradó feladatokat pedig olyan eszközökkel próbálják támogatni, melyek segítik az ember és a gép közötti interakciót (pl. releváns dokumentáció megjelenítésével, választási lehetőségek felajánlásával), illetve automatizáltan képesek ellenőrizni annak eredményét.

Ilyen komponenseket és eszközöket kínál például az Eclipse Foundation gondozásában álló Eclipse Modeling Project (\url{http://www.eclipse.org/modeling/}).
A projekt egyik legfontosabb és alapvető része az Eclipse Modeling Framework (röviden \gls{EMF}), mely egy, az iparban széles körben elterjedt, modell alapú szoftverfejlesztést támogató keretrendszer.
Az \gls{EMF} által nyújtott platformra épít az EMF-IncQuery keretrendszer, amely \gls{EMF} modellek felett teszi lehetővé gráfminta alapú lekérdezések deklaratív módon történő definiálását és hatékony végrehajtását.
(Az \gls{EMF} és EMF-IncQuery rendszereket a \ref{chap:techPrim}. fejezet részletesen ismerteti.)

Az EMF-IncQuery a lekérdezések specifikálására elsősorban egy szöveges leírónyelvet használ, melyre a későbbiekben legtöbbször \emph{lekérdező-} vagy \emph{tárgynyelvként} fogok hivatkozni az egyszerűség kedvéért.
A lekérdezések tárgynyelven történő megfogalmazása nagyon hasonló más programozási és leírónyelvek használatához.
Ez az analógia az elkövethető hibákban és -- szerencsére -- az ezek elleni védekezés módszereiben is visszaköszön.
Az egyik ilyen hibaelkerülési módszer, mellyel a legtöbb modern fordítóban és integrált fejlesztői környezetben is találkozhatunk, a \emph{kód statikus analízise}.

\section{Statikus analízis}

A \emph{statikus analízis} egy gyűjtőfogalom a fejlesztési időben rendelkezésre álló adatok alapján elvégezhető vizsgálatokra és információgyűjtésre, melyek a szintaxis kiemelésétől egészen az olyan formális módszerekig terjedhetnek, melyek a kód bizonyos szemantikai tulajdonságait matematikailag bizonyítják.
A spektrum egyik végén található szintaxis kiemelés a statikus analízis vizsgálatok egyik legegyszerűbb formája, mely szinte minden magára valamit adó szövegszerkesztőben megtalálható.
Egy fokkal kifinomultabbak, a már felületes szemantikai analízist is végező, úgynevezett ``lint'' eszközök, melyek a kódban gyakran előforduló, de formálisan nem vagy csak nehezen bizonyítható hibákra, hibalehetőségekre hívják fel a fejlesztők figyelmét.
A spektrum másik végén álló formális módszerekkel leggyakrabban valós idejű és biztonságkritikus rendszereknél (forgalom- és járműirányítás, nukleáris és orvosi rendszerek) találkozhatunk \cite{Storey:1996:SCC:524721}.
Ezeknél a rendszereknél emberi életek múlhatnak a szoftver helyes működésén, ezért az ellenőrzésekhez olyan technikákat használnak mint a modell ellenőrzés\footnote{A modell ellenőrzés -- a neve ellenére -- viszonylag kevéssé kötődik a dolgozat témájához.}, az adatfolyam analízis, az absztrakt interpretáció, a szimbolikus végrehajtás vagy a Hoare logika.

% a statikus analízis előnyei
A statikus analízis előnye, hogy relatíve hamar, még fejlesztési időben (pl. programírás közben, mentéskor vagy fordításkor) hívja fel a mérnökök figyelmét a problémákra, akik így időt takaríthatnak meg.
Egyrészt kevesebb automatizált tesztelési célt szolgáló teszteset megírására van szükség, hiszen a hibák egy részét a statikus elemzés segítségével már kizárták.
Másrészt a tesztelés -- akár automatizált, akár nem -- a gyakorlatban gyakran nem-determinisztikus, ellentétben a statikus analízis formális módszereivel.
Végül, de nem utolsó sorban, megspórolható a fejlesztő és a tesztelést végző személyzet -- rosszabb esetben a végfelhasználó -- közti kommunikációra fordított idő.

% a statikus analízis hátrányai
Ugyanakkor a statikus analízis sem nyújt univerzális megoldást minden problémára -- mivel általános esetben ez a megállási probléma megoldását jelentené, melyről tudjuk, hogy lehetetlen --, és sokszor csak ``lint'' típusú ellenőrzések futtathatóak a rendelkezésre álló korlátos idő és erőforrások miatt.
Ezek szükségszerűen konzervatív vizsgálatok, melyek nem mindig képesek hibák egyértelmű és biztos azonosítására -- sok a fals pozitív --, ezért kimenetük vizsgálata további munkát igényel.
Tehát a statikus analízis sem királyi út, érdemes futásidejű módszerekkel kombináltan használni.

Az eddig említésre került statikus vizsgálatok mind a fejlesztés segítésére vagy hibák, lehetséges hibaforrások felkutatására irányultak, még a fejlesztés befejezése előtt.
Ám a statikus analízis segítségével kinyert információk alapján javítható a futásidejű működés hatékonysága is, hiszen az elemzés során olyan optimalizációs lehetőségekre lelhetünk, melyeket beépíthetünk programunkba -- vagy éppen elhagyhatunk miatta részeket, ahogyan azt a nem elérhető kód elhagyásánál (dead code elimination) tesszük.

A statikus analízis azonban nem csak programokon értelmezhető fogalom, alkalmazható modellek vizsgálatára is, például előre nem ismert példánymodellek tulajdonságainak vizsgálatára metamodelljük alapján, hiszen a meta- és a példánymodell kapcsolata hasonló a forráskód és a futó program kapcsolatához.

\section{A dolgozat felépítése}

A \ref{chap:techPrim}. fejezetben először írok a modell-vezérelt szoftverfejlesztést támogató, és az EMF-IncQuery alapjául szolgáló Eclipse platformról és Eclipse Modeling Framework-ről.
Ezután bemutatom az EMF-IncQuery-t és példákon keresztül ismertetem a használatát és a lekérdezőnyelvét.

A \ref{chap:varUsage}. fejezetben olyan ellenőrzéseket tervezek és készítek el, melyek az EMF-IncQuery lekérdezőnyelvében található minták változóira történő hivatkozások minőségi és mennyiségi elemzése alapján hívják fel a fejlesztők figyelmét hibákra és hibalehetőségekre.

A \ref{chap:funcDep}. fejezetben a lekérdezőnyelvet a mintákban található változók között fennálló funkcionális függőségek szempontjából elemző ellenőrzéseket tervezek és készítek, melyek a kód és a kapcsolódó domén-modell között próbálnak ellentmondásokat felfedezni.
Továbbá egy példán keresztül röviden bemutatom, hogyan használható fel az általam előzőekben tárgyalt funkcionális függőségek elemzéséhez készített kód bizonyos lekérdezések végrehajtásának hatékonyabbá tételére.

Végül az \ref{chap:summary}. fejezetben összefoglalom a munkám eredményét és röviden beszélek az elkészült komponensek továbbfejlesztési lehetőségeiről.

%----------------------------------------------------------------------------
\chapter{Változó használat vizsgálat} %\label{sect:TechnologyPrimer}
%----------------------------------------------------------------------------
\section{Kutatás}
%----------------------------------------------------------------------------

%----------------------------------------------------------------------------
\chapter{Funkcionális függőség vizsgálat}
%----------------------------------------------------------------------------

%----------------------------------------------------------------------------
\section{Kutatás}
%----------------------------------------------------------------------------



%----------------------------------------------------------------------------
\chapter{Összefoglalás}
%----------------------------------------------------------------------------




\listoffigures\addcontentsline{toc}{chapter}{Ábrák jegyzéke}
%\listoftables\addcontentsline{toc}{chapter}{Táblázatok jegyzéke}

\bibliography{bibliography}\addcontentsline{toc}{chapter}{Irodalomjegyzék}
\bibliographystyle{plain}

\label{page:last}
\end{document}
